% !TEX TS-program = lualatex+se+debug
% !TEX encoding = UTF-8 Unicode

\documentclass[14pt, twoside, a4paper]{extreport}
% -----------------------------------------------------------------------------------
% ------------------------- USTAWIENIA STRONY--------------------------
% -----------------------------------------------------------------------------------
\usepackage{geometry}
\geometry{
	left=10mm,
	right=10mm,
	top=15mm,
	bottom=15mm,
	marginparwidth=0mm,
	marginparsep=0mm
	}
\newcommand{\wstawplik}[1]{\input{../include/#1}}%skraca kod
\wstawplik{headerofficium}
% elementy zmienne co do headerofficium:
\graphicspath{{../img/}} % katalog obrazów
%\svgpath{{../include/}}
\grechangestaffsize{20}%rozmiar wyniku GregorioTEX
% koniec 
\gresetlastline{ragged}% justowanie do lewej
% debugowanie:
% %\iffalse %*** %\fi %***
% \tracingall przed problematycznym kodem - bardzo dużo informacji
% lualatex+se+debug - dodatkowe opcje: -file-line-error i -halt-on-error 
% \errorcontextlines
%\def\mydebug{1}
\begin{document}
\makeatletter
\grechangedim{interwordspacetext}{\unexpanded{\the\fontdimen2\font plus \the\gluestretch\gre@space@skip@interwordspacenotes minus \the\glueshrink\gre@space@skip@interwordspacenotes}}{fixed}%
\makeatother
\pagestyle{empty}
%\pagestyle{plain}	
\large
\begin{spacing}{0.7}
\begin{center}
	\section*{In Parasceve - lectio 6}
	%\small\textit{Fontevraud Gradual\\}
	\small\textit{Laudes festivae (1940) p.84}
	\vspace{5mm}
\end{center}	
\end{spacing}
%\greannotation{} 
\gresetinitiallines{1}\greannotation{Lectio 6}\grecommentary{}\gabcsnippet{(f3)E(e)xa(f)cu(h)é(h)runt(h) tam(h)quam(h) glá(h)di(h)um(h) lin(h)guas(h) su(i)as.(f) (;) Non(h) di(h)cant(h) Iu(h)dǽ(h)i:(h) (,) Non(h) oc(h)cí(h)di(e)mus(f) Chri(g)stum.(f) (:) E(c)te(ef)nim(f) prop(f)té(f)re(f)a(f) e(f)um(f) de(f)dé(f)runt(f) iú(f)di(f)ci(g) Pi(h)lá(i)to,(h) (;) ut(h) qua(h)si(h) i(h)psi(h) a(h) mor(h)te(h) e(h)ius(h) vi(h)de(h)rén(h)tur(e) im(f)mú(g)nes.(f) (:) Nam(e) cum(f) di(h)xís(h)set(h) e(h)is(h) Pi(h)lá(i)tus:(f) (;) Vos(h) e(h)um(g) oc(f)cí(g)di(g)te:(g) (;) re(f)spon(g)dé(h)runt,(g) (;) No(i)bis(g) non(h) li(g)cet(f) (,) oc(f)cí(f)de(e)re(f) quem(g)quam.(f) (:) I(c)ni(e)qui(f)tá(f)tem(f) fa(f)cí(f)no(f)ris(e) su(f)i(f) (,) in(e) iú(f)di(h)cem(h) hó(h)mi(h)nem(h) re(h)fún(h)de(h)re(g) vo(f)lé(g)bant:(gh) (;) sed(h) num(h)quid(h) De(h)um(h) iú(h)di(h)cem(h) fal(f)lé(h)bant?(i) (:) Quod(h) fe(h)cit(h) Pi(h)lá(h)tus,(f) (;) in(h) e(h)o(h) i(h)pso(h) quod(h) fe(h)cit,(h) (,) a(h)li(h)quán(h)tum(h) pár(h)ti(g)ceps(f) fu(gh)it:(g) (;) sed(e) in(gh) com(f)pa(f)ra(f)ti(f)ó(f)ne(g) il(h)ló(i)rum(h) (,) mul(h)to(h) i(h)pse(h) in(e)no(f)cén(g)ti(f)or.(f) (:) In(c)sti(e)tit(f) e(f)nim(f) quan(f)tum(e) pó(f)tu(f)it,(f) (;) ut(e) il(f)lum(h) ex(h) e(h)ó(h)rum(h) má(h)ni(h)bus(h) li(g)be(f)rá(g)ret:(gh) (;) nam(h) prop(h)té(h)re(h)a(h) fla(h)gel(h)lá(h)tum(h) pro(h)dú(h)xit(e) ad(f) e(g)os.(f) (:) Non(c) per(ef)se(f)quén(fg)do(f) Dó(f)mi(f)num(f) fla(e)gel(f)lá(h)vit,(h) (;) sed(h) e(h)ó(h)rum(h) fu(h)ró(h)ri(h) sa(h)tis(h)fá(h)ce(g)re(f) vo(gh)lens:(g) (;) ut(e) vel(gh) sic(f) iam(f) mi(e)té(f)sce(f)rent,(f) (,) et(f) de(f)sí(f)ne(f)rent(f) vel(f)le(g) oc(h)cí(i)de(h)re,(h) (;) cum(h) fla(h)gel(h)lá(h)tum(e) vi(f)dé(g)rent.(f) (:) Fe(f)cit(e) et(c) hoc.(ef) (:) At(e) u(f)bi(h) per(h)se(h)ve(h)ra(h)vé(i)runt,(f) (;) no(h)stis(h) il(h)lum(h) la(h)vís(h)se(h) ma(h)nus,(h) et(g) di(f)xís(g)se,(gh) (;) quod(h) i(h)pse(h) non(h) fe(h)cís(h)set,(h) (,) mun(h)dum(h) se(h) es(h)se(h) a(h) mor(h)te(e) il(f)lí(g)us.(f) (:) Fe(f)cit(ec) ta(e)men.(ef) (:) Sed(e) si(f) re(h)us,(h) qui(h)a(h) fe(h)cit(h) vel(g) in(f)ví(gh)tus:(g) (;) il(e)li(gh) in(f)no(e)cén(f)tes,(f) (,) qui(f) co(f)ë(f)gé(f)runt(g) ut(h) fá(i)ce(i)ret?(i) (:) Nul(h)lo(ef) mo(g)do.(f) (:) Sed(c) il(e)le(f) di(f)xit(f) in(f) e(f)um(f) sen(e)tén(f)ti(f)am,(f) (,) et(f) ius(f)sit(f) e(f)um(f) cru(g)ci(h)fí(i)gi,(h) (;) et(h) qua(h)si(h) i(h)pse(g) oc(f)cí(gh)dit:(g) (;) et(e) vos,(gh) o(f) Iu(f)dǽ(f)i,(f) oc(e)ci(f)dí(g)stis.(f) (:) Un(e)de(f) oc(h)ci(f)dí(h)stis?(i) (:) Glá(h)di(g)o(f) lin(g)guæ:(gh) (;) a(h)cu(h)í(h)stis(h) e(h)nim(h) lin(e)guas(f) ve(g)stras.(f) (:) Et(e) quan(f)do(h) per(h)cus(h)sí(h)stis,(h) (,) ni(h)si(h) quan(h)do(h) cla(h)má(i)stis:(f) (;) Cru(h)ci(h)fí(h)ge,(f) (,) cru(h)ci(h)fí(h)ge?(hi) (::)}\minispace%
\tlumacz{Wyostrzyli języki swoje jak miecz. Niech nie mówią Żydzi: Nie myśmy zabili Chrystusa. Gdyż jeśli go zaprowadzili przed trybunał Piłata, to na to, ażeby wydać się niejako niewinnymi Jego śmierci. I w istocie, gdy Piłat im powiedział: Sami skażcie go, na śmierć, odpowiedzieli: Niewolno nam nikogo ua śmierć skazywać. Chcieli więc, by cała odpowiedzialność za zbrodnię spadła na sędziego człowieka; lecz zaliż mogli oszukać Sędziego Boga? W tem, co uczynił Piłat, stał się po części uczestnikiem zbrodni, lecz w porówaniu z nimi, nierównie mniej jest winnym: gdyż nalegał ile tylko mógł, aby go z rąk ich wyrwać: i przeto okrutnie ubiczowanego przed nimi stawił. Nie z chęci prześladowania ubiczował Pana, lecz zajadłości ich chciał dogodzić, aby się tem przynajmniej ułagodzili i zaniechali zamiaru zabicia go, gdy tak dalece ubiczowanego ujrzą. Uczynił i to. Gdy zaś uporczywie trwali w złości, wiecie, iż umył ręce swoje i oświadczył się z tem, że nie on to popełnił; że nie winien śmierci jego. Jednakże uczynił. Jeżeli więc on jest winien, chociaż niechętnie zdał go na śmierć, mająż być niewinnymi, którzy go zmusili, aby to uczynił? Żadną miarą. Lecz on wydał na niego wyrok i rozkazał, aby go ukrzyżowano, a już tem samem jakby go zabił; z tem wszystkiem wy, o Żydzi, zabiliście. Jakimże sposobem zabiliście? Oto mieczem języka, bo zaostrzyliście języki wasze. Kiedyż mianowicie zadaliście cios śmiertelny, jeśli nie wtedy, gdyście krzyczeli: Ukrzyżuj, ukrzyżuj.}
\end{document}